\documentclass[11pt]{article}
\usepackage[margin=1in]{geometry}
\usepackage{amsmath, amssymb, amsthm, amsfonts}
\usepackage{booktabs}
\usepackage{algorithm}
\usepackage{algorithmic}
\usepackage{enumitem}
\usepackage{hyperref}
\usepackage{xcolor}
\usepackage{bm}
\usepackage{tikz}

\renewcommand{\algorithmicrequire}{\textbf{Input:}}
\renewcommand{\algorithmicensure}{\textbf{Output:}}

\DeclareMathOperator*{\argmin}{arg\,min}
\DeclareMathOperator*{\argmax}{arg\,max}
\DeclareMathOperator{\Flow}{Flow}
\DeclareMathOperator{\Suff}{Suff}
\DeclareMathOperator{\Comp}{Comp}
\DeclareMathOperator{\Mono}{Mono}
\DeclareMathOperator{\FaithCorr}{FaithCorr}
\DeclareMathOperator{\layer}{layer}
\DeclareMathOperator{\sign}{sign}
\DeclareMathOperator{\SCS}{SCS}

\newtheorem{definition}{Definition}
\newtheorem{proposition}{Proposition}
\newtheorem{lemma}{Lemma}

\title{Flow-Based XAI Evaluation for \\ Attribution Graph Summarization}
\author{}
\date{}

\begin{document}
\maketitle

\section{Motivation}\label{sec:motivation}

A standard requirement for any attribution-based explanation is the
\emph{completeness axiom}: the sum of the parts equals the whole.
A naive operationalization---summing supernode-to-logit edge
weights and comparing against the model's logit response---fails
because the attribution graph and the model's logits live in
different units. Attribution edge weights $W$ are derived from a
local linear surrogate (the cross-layer transcoder); logits are
nonlinear forward-pass outputs. There is no theoretical guarantee
that $\sum_{S} \mathcal{A}(S \to V_{\text{logit}})$ matches
$L(x) - L_0$, so any measure built from that comparison is
ill-posed.

We resolve this by defining all evaluation quantities in a
\emph{single unit}: graph flow. The original attribution graph
$G = (V, E, W)$ is a directed acyclic graph with a distinguished
source set (embedding nodes $V_{\text{emb}}$) and sink set (logit
nodes $V_{\text{logit}}$). On such a graph there is a natural
notion of flow from sources to sinks; it satisfies a conservation
law at every interior node by construction. Faithfulness
(``important parts cause large drops when removed'') and
completeness (``the parts sum to the whole'') then become two
facets of the same flow quantity, with no unit mismatch.

This document (i) defines flow on the original attribution graph,
(ii) lifts it to the supernode graph, (iii) defines completeness,
deletion, insertion, and ranking measures in flow units, and
(iv) gives the computational procedure.

\section{Flow on the Original Graph}\label{sec:flow}

\subsection{Setup}

The attribution graph $G = (V, E, W)$ is a DAG with vertex partition
$V = V_{\text{emb}} \cup V_{\text{feat}} \cup V_{\text{logit}}$.
Embedding nodes have no in-edges; logit nodes have no out-edges.
Each node carries a layer index $\layer(\cdot)$ with embeddings at
layer 0 and logits at the final layer. The weighted adjacency
$W \in \mathbb{R}^{N \times N}$ has $W_{uv}$ the signed attribution
weight of edge $(u, v)$.

\subsection{Source flow and conservation}

We treat the embedding nodes as flow sources and propagate flow
forward through the DAG. To handle signed weights cleanly we work
with absolute weights $|W|$ throughout, treating sign as a separate
attribute used only in the directional measures of
Section~\ref{sec:directional}.

\paragraph{Outgoing weight normalization.}
For each non-sink node $u$, define the outgoing transition matrix
\begin{equation}
P^{\text{out}}_{uv} = \frac{|W_{uv}|}{\sum_{w : (u, w) \in E} |W_{uw}|},
\label{eq:Pout}
\end{equation}
with $\sum_v P^{\text{out}}_{uv} = 1$ whenever $u$ has at least one
out-edge. This redistributes any flow entering $u$ across its
out-edges in proportion to attribution mass.

\paragraph{Source vector.}
Let $\tau \in \mathbb{R}^{|V_{\text{emb}}|}$ be the token-importance
vector of the main paper (normalized SHAP values),
$\sum_{e \in V_{\text{emb}}} \tau_e = 1$. We inject one unit of
flow into the embedding layer, partitioned by $\tau$:
\begin{equation}
F^{\text{src}}_u =
\begin{cases}
\tau_u & u \in V_{\text{emb}}, \\
0      & \text{otherwise}.
\end{cases}
\label{eq:Fsrc}
\end{equation}

\paragraph{Node flow.}
The flow at node $u$ is the total mass arriving via all directed
source-to-$u$ paths:
\begin{equation}
\Flow(u) = \bigl[\, F^{\text{src}\top}\,(I - P^{\text{out}})^{-1}\,\bigr]_u,
\label{eq:flow}
\end{equation}
which converges on the DAG because $P^{\text{out}}$ is strictly
upper-triangular in any topological ordering and so nilpotent.
Equivalently, $\Flow$ satisfies the linear recursion
\begin{equation}
\Flow(u) = F^{\text{src}}_u + \sum_{w : (w, u) \in E} \Flow(w) \cdot P^{\text{out}}_{wu}.
\label{eq:flow-recursion}
\end{equation}

\paragraph{Edge flow.}
The flow on edge $(u, v)$ is the mass that traverses it:
\begin{equation}
\Flow(u, v) = \Flow(u) \cdot P^{\text{out}}_{uv}.
\label{eq:edge-flow}
\end{equation}

\subsection{Conservation lemma}

\begin{lemma}[Flow conservation]\label{lem:conservation}
For any node $u \notin V_{\text{emb}} \cup V_{\text{logit}}$,
\begin{equation}
\sum_{w : (w, u) \in E} \Flow(w, u)
\;=\;
\sum_{v : (u, v) \in E} \Flow(u, v)
\;=\;
\Flow(u).
\label{eq:conservation}
\end{equation}
Furthermore, total flow is preserved end-to-end:
\begin{equation}
\sum_{\ell \in V_{\text{logit}}} \Flow(\ell)
\;=\;
\sum_{e \in V_{\text{emb}}} F^{\text{src}}_e
\;=\;
1.
\label{eq:total-flow}
\end{equation}
\end{lemma}

\begin{proof}[Proof sketch]
The recursion~\eqref{eq:flow-recursion} together with
$\sum_v P^{\text{out}}_{uv} = 1$ for interior nodes gives the per-node
balance~\eqref{eq:conservation}. Summing the balance over all
interior nodes telescopes along directed paths, leaving the
source-sink identity~\eqref{eq:total-flow}.
\end{proof}

Equation~\eqref{eq:total-flow} is the key property: total flow is
$1$ by construction, regardless of the graph, the partition, or
the attribution weights. Completeness is automatic on the original
graph; the question becomes how much of it survives pruning and
supernode aggregation.

\section{Flow on the Supernode Graph}\label{sec:supernode-flow}

\subsection{Lifting flow to supernodes}

Given a supernode partition $\varphi$ of the pruned graph
$(V', E')$, contract each supernode $S \in \pi(\varphi)$ into a
single vertex. Embedding and logit nodes remain singletons by
construction. The supernode flow is defined by aggregating the
node flow:
\begin{equation}
\Flow^{SN}(S) = \sum_{u \in S} \Flow(u),
\label{eq:flow-sn-node}
\end{equation}
and supernode edge flow by aggregating edge flows of original edges
crossing the cluster boundary:
\begin{equation}
\Flow^{SN}(S, T) = \sum_{\substack{(u, v) \in E' \\ u \in S, v \in T}} \Flow(u, v).
\label{eq:flow-sn-edge}
\end{equation}

\begin{lemma}[Supernode conservation]\label{lem:sn-conservation}
For any non-source, non-sink supernode $S$,
$\sum_{R} \Flow^{SN}(R, S) = \sum_{T} \Flow^{SN}(S, T) = \Flow^{SN}(S)$.
\end{lemma}

The proof is immediate: summing~\eqref{eq:conservation} over
$u \in S$ and splitting flows by whether their endpoints lie in
$S$ or its complement gives the supernode balance.

\subsection{Pruning and flow loss}

Pruning $(V, E) \to (V', E')$ in Block~1 of the main paper deletes
nodes and edges with low combined-evidence score. Let
$\Flow_{V'}(\cdot)$ denote the flow computed on the pruned graph
(with the same source vector $\tau$ restricted to surviving
embedding nodes, renormalized to sum to 1). The
\emph{pruning flow loss} is
\begin{equation}
\mathcal{L}_{\text{prune}} = 1 - \sum_{\ell \in V_{\text{logit}}}
\Flow_{V'}(\ell).
\label{eq:prune-loss}
\end{equation}
By Lemma~\ref{lem:conservation} this is exactly the flow that
\emph{would have} reached the logits via paths through pruned nodes
or edges. The pruning thresholds $(\eta_V, \eta_E)$ in the main
paper are chosen so that $\mathcal{L}_{\text{prune}}$ is small;
$1 - \mathcal{L}_{\text{prune}}$ is the fraction of flow retained
and provides a natural sanity check on Block~1.

\section{Flow-Based Completeness}\label{sec:completeness}

\begin{definition}[Flow Completeness]\label{def:flow-completeness}
For a supernode partition $\varphi$ on the pruned graph,
\begin{equation}
\mathcal{C}(\varphi) = \sum_{\ell \in V_{\text{logit}}} \Flow^{SN}(\ell).
\label{eq:flow-completeness}
\end{equation}
\end{definition}

\begin{proposition}[Bounds and interpretation]
$\mathcal{C}(\varphi) \in [0, 1]$, and
\begin{equation}
\mathcal{C}(\varphi) = 1 - \mathcal{L}_{\text{prune}}.
\label{eq:completeness-equals-retention}
\end{equation}
In particular, $\mathcal{C}(\varphi)$ does not depend on the
partition $\varphi$ (only on the pruned graph $(V', E')$): supernode
aggregation conserves flow.
\end{proposition}

\begin{proof}
Singleton logit supernodes give
$\Flow^{SN}(\ell) = \Flow_{V'}(\ell)$, so the sum
in~\eqref{eq:flow-completeness} is exactly the pruned-graph total
flow into logits, which equals $1 - \mathcal{L}_{\text{prune}}$ by
Lemma~\ref{lem:conservation} applied to the pruned graph.
\end{proof}

This is the right shape for a completeness axiom: well-defined,
bounded in $[0, 1]$, equal to 1 when no flow is lost, and invariant
under valid contractions. It also makes precise what $\SCS$ was
trying to capture: \emph{completeness measures pruning retention,
not partition quality}. All clustering methods that share the same
Block~1 output report identical $\mathcal{C}(\varphi)$.

\section{Flow-Based Faithfulness via Supernode Deletion}\label{sec:faithfulness}

Faithfulness asks whether \emph{removing} important supernodes
costs more flow than removing unimportant ones. Define supernode
deletion as removing $S$ from the pruned graph along with all
edges incident to it, then recomputing flow on the surviving
subgraph.

\subsection{Deletion flow and supernode importance}

\begin{definition}[Deletion flow]\label{def:deletion-flow}
Let $G_{SN} \setminus S$ denote the supernode graph with $S$ and
its incident edges removed, and let $\Flow_{\setminus S}$ be the
flow on this subgraph (with the same source vector). The
\emph{deletion flow loss} of $S$ is
\begin{equation}
\Delta \mathcal{C}(S)
= \mathcal{C}(\varphi)
- \sum_{\ell \in V_{\text{logit}}} \Flow_{\setminus S}(\ell).
\label{eq:deletion-flow}
\end{equation}
\end{definition}

$\Delta \mathcal{C}(S) \in [0, \mathcal{C}(\varphi)]$ measures the
fraction of logit-bound flow that passes through $S$. It is the
flow-theoretic analog of an ablation effect: deleting a supernode
that lies on heavy source-to-sink paths produces a large
$\Delta \mathcal{C}$, while deleting a side-branch supernode
produces a small one.

The natural importance score is the supernode's own flow:
\begin{equation}
\sigma_{\text{flow}}(S) = \Flow^{SN}(S).
\label{eq:sigma-flow}
\end{equation}

\subsection{Per-supernode faithfulness correlation}

\begin{definition}[Flow Faithfulness Correlation]\label{def:flow-faithcorr}
\begin{equation}
\FaithCorr_{\Flow}(\varphi)
= \mathrm{Spearman}\bigl(
    \sigma_{\text{flow}}(S),\;
    \Delta \mathcal{C}(S)
  \bigr)_{S \in \pi(\varphi)}.
\label{eq:flow-faithcorr}
\end{equation}
A faithful partition ranks supernodes by flow consistently with
their causal contribution to logit-bound flow.
\end{definition}

A perfectly tree-shaped flow graph yields
$\FaithCorr_{\Flow} = 1$: each supernode's deletion loss equals its
own flow exactly. Divergence from 1 measures how much
\emph{flow redundancy} the partition contains: alternative paths
that route around a deleted supernode and preserve sink flow.

\subsection{Deletion AUC}

\begin{definition}[Flow Deletion AUC]\label{def:flow-deletion-auc}
Order supernodes by decreasing $\sigma_{\text{flow}}$:
$S_{(1)}, \ldots, S_{(|\pi|)}$. Let $\mathcal{S}_r = \{S_{(1)}, \ldots, S_{(r)}\}$.
Define the deletion curve
\begin{equation}
\delta(r) = \frac{1}{\mathcal{C}(\varphi)} \sum_{\ell \in V_{\text{logit}}}
\Flow_{\setminus \mathcal{S}_r}(\ell),
\end{equation}
and the deletion AUC
\begin{equation}
\Flow\text{-}\mathrm{DAUC}(\varphi)
= \frac{1}{|\pi(\varphi)| + 1} \sum_{r=0}^{|\pi(\varphi)|} \delta(r).
\label{eq:flow-dauc}
\end{equation}
$\delta(0) = 1$, $\delta(|\pi(\varphi)|) = 0$;
\emph{lower $\Flow\text{-}\mathrm{DAUC}$ is more faithful}.
\end{definition}

\subsection{Insertion AUC}

\begin{definition}[Flow Insertion AUC]\label{def:flow-insertion-auc}
Let $G_{SN}\!\restriction_{\mathcal{S}_r}$ denote the subgraph
keeping only supernodes in $\mathcal{S}_r$ together with the
embedding and logit singletons. The insertion curve is
\begin{equation}
\iota(r) = \frac{1}{\mathcal{C}(\varphi)} \sum_{\ell \in V_{\text{logit}}}
\Flow_{G_{SN}\restriction \mathcal{S}_r}(\ell),
\end{equation}
and
\begin{equation}
\Flow\text{-}\mathrm{IAUC}(\varphi)
= \frac{1}{|\pi(\varphi)| + 1} \sum_{r=0}^{|\pi(\varphi)|} \iota(r).
\label{eq:flow-iauc}
\end{equation}
\emph{Higher is more faithful.}
\end{definition}

\subsection{Sufficiency and Comprehensiveness}

\begin{definition}[Flow Sufficiency at top-$k$]\label{def:flow-suff}
\begin{equation}
\Suff_k^{\Flow}(\varphi) = 1 - \iota(k).
\label{eq:flow-suff}
\end{equation}
Lower is better: the top-$k$ supernodes alone should retain most of
the sink flow.
\end{definition}

\begin{definition}[Flow Comprehensiveness at top-$k$]\label{def:flow-comp}
\begin{equation}
\Comp_k^{\Flow}(\varphi) = 1 - \delta(k).
\label{eq:flow-comp}
\end{equation}
Higher is better: removing the top-$k$ supernodes should cost most
of the sink flow.
\end{definition}

\subsection{Monotonicity}

\begin{definition}[Monotonicity]\label{def:mono}
With $\Delta \mathcal{C}_{(k)} = \Delta \mathcal{C}(S_{(k)})$,
\begin{equation}
\Mono(\varphi)
= \frac{1}{|\pi(\varphi)| - 1}
  \sum_{k=1}^{|\pi(\varphi)| - 1}
  \mathbf{1}\bigl[\, \Delta \mathcal{C}_{(k)} \ge \Delta \mathcal{C}_{(k+1)} \,\bigr].
\label{eq:mono}
\end{equation}
Higher is more faithful.
\end{definition}

\section{Directional Faithfulness}\label{sec:directional}

Flow as defined above uses $|W|$ and is sign-agnostic. To recover
directional information, we track the signed contribution
separately. Let
\begin{equation}
\Flow^{\pm}(u, v) = \Flow(u, v) \cdot \sign(W_{uv})
\label{eq:flow-signed}
\end{equation}
and aggregate at the supernode level by net inflow:
\begin{equation}
\Flow^{SN, \pm}(S) = \sum_{e \in V_{\text{emb}}} \sum_{\text{paths } e \to S}
\prod_{(u, v) \in \text{path}} \sign(W_{uv}) \cdot \Flow(u, v).
\end{equation}
A partition is \emph{directionally faithful} if features sharing a
supernode contribute with consistent sign. The opposing-sign
fraction of the main paper is exactly the empirical version of this
constraint at the pair level.

\section{Computational Procedure}\label{sec:procedure}

\begin{algorithm}[h]
\caption{Flow-based XAI evaluation of a supernode partition}
\label{alg:flow-xai}
\begin{algorithmic}[1]
\REQUIRE Pruned graph $(V', E', W')$; partition $\varphi$; source $\tau$;
top-$k$ values $\{k_1, \ldots, k_m\}$
\ENSURE Measures $(\mathcal{C}, \FaithCorr_{\Flow}, \Flow\text{-}\mathrm{DAUC},
                   \Flow\text{-}\mathrm{IAUC}, \Mono, \{\Suff_{k_j}^{\Flow}, \Comp_{k_j}^{\Flow}\})$
\STATE Compute $P^{\text{out}}$ via Eq.~\eqref{eq:Pout}
\STATE Compute $\Flow(u)$ for all $u \in V'$ via the
       recursion~\eqref{eq:flow-recursion} in topological order
       \COMMENT{$O(|E'|)$}
\STATE Compute $\Flow^{SN}(S)$ and $\Flow^{SN}(S, T)$ via
       Eqs.~\eqref{eq:flow-sn-node}--\eqref{eq:flow-sn-edge}
\STATE $\mathcal{C}(\varphi) \gets \sum_{\ell} \Flow^{SN}(\ell)$
\FOR{each $S \in \pi(\varphi)$}
  \STATE Recompute flow on $G_{SN} \setminus S$ via
         Eq.~\eqref{eq:flow-recursion} restricted to surviving nodes
         \COMMENT{$O(|E'|)$ per supernode}
  \STATE $\Delta \mathcal{C}(S) \gets$ Eq.~\eqref{eq:deletion-flow}
\ENDFOR
\STATE Sort supernodes by $\sigma_{\text{flow}} = \Flow^{SN}$
\FOR{$r = 0, 1, \ldots, |\pi(\varphi)|$}
  \STATE Compute $\delta(r)$ and $\iota(r)$ by recomputing flow on
         the corresponding subgraphs
\ENDFOR
\STATE Compute $\FaithCorr_{\Flow}, \Flow\text{-}\mathrm{DAUC},
       \Flow\text{-}\mathrm{IAUC}, \Mono$ from the curves
\STATE Compute $\Suff_{k_j}^{\Flow}, \Comp_{k_j}^{\Flow}$ for each $k_j$
\RETURN all measures
\end{algorithmic}
\end{algorithm}

\paragraph{Cost.}
Each flow computation is $O(|E'|)$ via topological-order
propagation. The total cost per partition is $O(|\pi(\varphi)| \cdot |E'|)$:
one base flow plus one flow per supernode for the deletion curve
(insertion reuses values via $\iota(r) = 1 - \delta(|\pi| - r)$ when
the graph factorizes; in general both directions are computed).
Crucially, \textbf{no forward passes through the model are
required}: the entire evaluation runs on the attribution graph
alone, using only its weights $W'$ and the source vector $\tau$.
This is a substantial cost reduction compared with the ablation-based
measures of the previous version.

\section{Why Flow Is the Right Unit}\label{sec:discussion}

\paragraph{Unit consistency.}
Completeness, deletion loss, sufficiency, and comprehensiveness are
all measured as fractions of total source-to-sink flow. There is no
comparison between attribution-derived quantities and forward-pass
logits, so no unit mismatch can arise.

\paragraph{Conservation by construction.}
Lemmas~\ref{lem:conservation} and~\ref{lem:sn-conservation} pin the
completeness axiom to the graph structure. Completeness fails only
when flow is destroyed---by pruning, by removing supernodes during
deletion tests, or by restricting to top-$k$ for insertion. Each
failure has a precise cause traceable to a graph operation.

\paragraph{Relationship to the main paper's Inf and Rel.}
The Inf and Rel quantities of Appendix~A.1 are special cases of
flow. Inf is the flow when $\tau$ is uniform and the propagation
direction is reversed (sink-to-source); Rel is forward flow with
$\tau$ as defined here. Edge-level Inf$_{uv}$ and Rel$_{uv}$ are
edge flows in their respective directions. The flow framework
unifies these into one quantity used consistently across pruning
(Block~1), clustering (Block~2 via $\tilde S$), and evaluation
(this document).

\paragraph{Faithfulness vs.\ flow path redundancy.}
$\FaithCorr_{\Flow} < 1$ has a sharp interpretation: it measures how
many supernodes lie on parallel source-to-sink paths whose removal
can be compensated by alternative routes. A high
$\FaithCorr_{\Flow}$ identifies a partition with low path
redundancy---each supernode is a bottleneck for the flow attributed
to it. This is exactly the property the supernode graph should have
to support causal reading.

\paragraph{Comparison to ablation-based faithfulness.}
Forward-pass ablation faithfulness (the earlier version of this
document) and flow-based faithfulness will agree on partitions
where the transcoder linearization is faithful to the model.
Reporting both is informative: large gaps diagnose transcoder
failure, not partition failure.

\section{Proposed Comparison Table}\label{sec:proposed-table}

\begin{table}[h]
\centering
\caption{Flow-based XAI evaluation on the analogies dataset
($N = 55$ pruned graphs, $K = \lfloor n/2 \rfloor$). Cost-style
($\downarrow$): $\Flow$-DAUC, $\Suff^{\Flow}$. Benefit-style
($\uparrow$): $\mathcal{C}$, $\FaithCorr_{\Flow}$,
$\Flow$-IAUC, $\Mono$, $\Comp^{\Flow}$.
$\mathcal{C}$ is partition-invariant and reported once per Block~1
configuration.}
\label{tab:flow-xai}
\small
\begin{tabular}{@{}lccccccc@{}}
\toprule
\textbf{Method} &
$\mathcal{C} \uparrow$ &
$\FaithCorr_{\Flow} \uparrow$ &
$\Flow$-DAUC $\downarrow$ &
$\Flow$-IAUC $\uparrow$ &
$\Mono \uparrow$ &
$\Suff_5^{\Flow} \downarrow$ &
$\Comp_5^{\Flow} \uparrow$ \\
\midrule
\textsc{ASG-arith} & --- & --- & --- & --- & --- & --- & --- \\
\textsc{ASG-harm}  & --- & --- & --- & --- & --- & --- & --- \\
\textsc{ASG-geo}   & --- & --- & --- & --- & --- & --- & --- \\
Modularity         & --- & --- & --- & --- & --- & --- & --- \\
Spectral-RBF       & --- & --- & --- & --- & --- & --- & --- \\
K-means            & --- & --- & --- & --- & --- & --- & --- \\
\bottomrule
\end{tabular}
\end{table}

\section{Limitations}\label{sec:limitations}

\begin{enumerate}[leftmargin=*,topsep=2pt,itemsep=1pt]
\item Flow uses $|W|$ and so loses sign information. Directional
faithfulness (Section~\ref{sec:directional}) recovers it but at
the cost of leaving the conservation framework.
\item Flow is a property of the attribution graph, not the model.
Faithfulness of the attribution graph itself (transcoder fidelity)
remains a separate question. Forward-pass ablation provides the
complementary test.
\item Deletion flow is the absolute drop in sink flow when $S$ is
removed; alternative definitions (relative drop, weighted by which
sinks lose flow) are possible and may better match
application-specific notions of importance.
\item The conservation law holds exactly only when
$P^{\text{out}}$ is row-stochastic at every interior node. Nodes
with no out-edges in the pruned graph violate this; we treat them
as sinks with zero flow forwarded, which is conservative for
completeness but inflates $\mathcal{L}_{\text{prune}}$.
\end{enumerate}

\end{document}